\section{Суперкомпьютерные технологии}

В принятой в пособии СПбПУ трактовке \textbf{суперкомпьютерные технологии} включают три взаимосвязанных блока: технологии создания высокопроизводительных вычислительных систем, системное и прикладное программное обеспечение и технологии их проблемно-ориентированного применения. Их объединяет использование параллельных вычислений для ресурсоёмких задач математического моделирования.

\subsection{Архитектура параллельных систем}

Современный суперкомпьютер обычно состоит из множества вычислительных узлов, соединённых высокоскоростной сетью. Узел содержит многоядерные CPU, оперативную память и, возможно, ускорители GPU.

По организации памяти принципиальны две модели:
\begin{itemize}
\item \textbf{разделяемая память}: потоки работают в общем адресном пространстве;
\item \textbf{распределённая память}: каждый процесс/узел имеет свою память, а данные передаются сообщениями.
\end{itemize}
Современные кластеры обычно гибридны: общая память внутри узла и распределённая между узлами.

Для производительности существенны не только число процессоров, но и иерархия памяти, пропускная способность, латентность сети и локальность данных.

\subsection{Параллельный алгоритм и декомпозиция}

Параллельный расчёт требует разбиения задачи на подзадачи. При \textbf{декомпозиции по данным} область или массив делится между процессами:
\[
\Omega=\bigcup_{k=1}^{p}\Omega_k.
\]
При функциональной декомпозиции различным вычислителям назначаются разные стадии алгоритма.

Хорошая декомпозиция должна обеспечивать:
\begin{itemize}
\item баланс вычислительной нагрузки;
\item малый объём обменов;
\item достаточно крупные независимые блоки вычислений;
\item возможность совмещения вычислений и коммуникаций.
\end{itemize}
Если один процесс выполняет существенно больше работы остальных, время всего расчёта определяется этим процессом.

\subsection{MPI и общая память}

Для распределённой памяти основным стандартом является \textbf{MPI} (Message Passing Interface). Параллельная программа запускается как несколько процессов с собственными адресными пространствами. Используются двухточечные обмены и коллективные операции --- рассылка, сбор, редукция и др.

Для суммы
\[
S=\sum_{i=1}^{N}a_i
\]
каждый процесс может вычислить локальную сумму $S_k$, после чего коллективная редукция даёт
\[
S=\sum_{k=1}^{p}S_k.
\]
Время сообщения грубо представляют моделью
\[
T_{\rm comm}\approx \alpha+\beta m,
\]
где $\alpha$ --- латентность, $m$ --- объём сообщения, а $\beta$ характеризует стоимость передачи единицы данных.

Для общей памяти применяются потоковые модели, например OpenMP: несколько потоков используют общее адресное пространство. Здесь возникают задачи синхронизации, гонки данных и издержки на критические секции.

В гетерогенных системах часть массово-параллельных вычислений может выполняться на GPU. Для больших кластеров типична гибридная схема
\[
\text{MPI между узлами}+\text{потоки или GPU внутри узла}.
\]

\subsection{Производительность и масштабируемость}

Пусть $T_1$ --- время последовательного решения, $T_p$ --- время на $p$ вычислителях. \textbf{Ускорение} и \textbf{эффективность} определяются как
\[
S_p=\frac{T_1}{T_p},
\qquad
E_p=\frac{S_p}{p}.
\]
Идеально $S_p=p$ и $E_p=1$, но на практике эффективность уменьшают последовательные участки, обмены, синхронизации и дисбаланс.

\textbf{Закон Амдала.} Если доля вычислений, допускающая идеальное распараллеливание, равна $\alpha$, то для задачи фиксированного размера
\[
S_p\le
\frac{1}{(1-\alpha)+\alpha/p},
\qquad
S_{\infty}\le \frac{1}{1-\alpha}.
\]
Следовательно, даже небольшой последовательный участок ограничивает максимальное ускорение.

Различают:
\begin{itemize}
\item \textbf{сильную масштабируемость}: размер задачи фиксирован, число процессоров увеличивается;
\item \textbf{слабую масштабируемость}: объём задачи растёт примерно пропорционально числу процессоров, так что работа на один процесс остаётся приблизительно постоянной.
\end{itemize}
В пособии СПбПУ масштабируемость рассматривается и как свойство вычислительной системы, и как способность параллельной программы эффективно использовать наращиваемые ресурсы.

\subsection{Память и арифметическая интенсивность}

Производительность может ограничиваться не числом арифметических устройств, а передачей данных из памяти. Полезной характеристикой является \textbf{арифметическая интенсивность}
\[
I=\frac{\text{число арифметических операций}}{\text{объём переданных данных}}.
\]
Алгоритмы с малой $I$ часто ограничены пропускной способностью памяти; с большой $I$ --- пиковой вычислительной производительностью. Поэтому в HPC важны локальность данных и уменьшение обменов между уровнями памяти и узлами.


\subsection{Закон Густафсона и измерение масштабируемости}

Закон Амдала отвечает на вопрос о фиксированной задаче. На практике с ростом
числа вычислителей часто увеличивают и размер задачи. Если $s$ --- доля
последовательного времени в уже масштабированной задаче, то закон Густафсона
оценивает ускорение как
\[
\boxed{S_p^{(G)}=p-s(p-1)}.
\]
Тем самым небольшая последовательная часть не мешает эффективно использовать
большое число процессоров, если параллельная часть задачи растёт вместе с $p$.

При экспериментальном исследовании масштабируемости полезно строить не только
$S_p$, но и $E_p$, время коммуникаций и долю дисбаланса. Для сильного
масштабирования фиксируют размер задачи; для слабого --- сохраняют примерно
постоянный объём работы на процесс. Это позволяет отличить ограничение самого
алгоритма от недостатка параллельной работы на выбранном размере задачи.

\subsection{Коллективные операции и коммуникационная сложность}

Для распределённых алгоритмов важен не только общий объём пересылаемых данных,
но и число синхронизаций. Стоимость коллективной редукции, рассылки или
all-to-all зависит от топологии сети и реализации MPI; поэтому операция,
выполняемая на каждом шаге временного интегрирования, способна ограничить
масштабируемость даже при небольшом размере сообщения.

Полезный принцип --- укрупнять сообщения, уменьшать число глобальных барьеров и,
где возможно, перекрывать неблокирующие обмены полезными вычислениями. Для
сеточных задач также стремятся уменьшить отношение размера межпроцессорной
границы к объёму локальной подобласти.

\subsection{Roofline-оценка и локальность данных}

Связь вычислительной производительности с уже введённой арифметической
интенсивностью удобно выражать простейшей roofline-оценкой
\[
\boxed{P\le \min(P_{peak},\;B_{mem}I)},
\]
где $P_{peak}$ --- пиковая производительность вычислительных устройств,
$B_{mem}$ --- пропускная способность памяти. Она сразу показывает два режима:
memory-bound и compute-bound.

Отсюда следует практическая важность блочных алгоритмов, последовательного
доступа, повторного использования данных из кэша и уменьшения лишних копирований.
Оптимизация числа FLOPS сама по себе может почти не ускорить memory-bound ядро.

\subsection{Гибридные вычисления}

На современных кластерах естественна гибридная схема: MPI распределяет работу
между узлами, OpenMP или потоки используют ядра общей памяти внутри узла, а GPU
выполняет массово-параллельные ядра. Эффективность требует согласованной
декомпозиции, минимизации копирований CPU--GPU и перекрытия коммуникаций
вычислениями.

Для больших расчётов также важны контрольные точки (checkpoint/restart),
детерминированное хранение параметров запуска и мониторинг производительности:
сбой одного длинного запуска не должен приводить к полной потере многосуточного
расчёта.


\subsection{Что важно сказать на экзамене}

Главная цепочка ответа:
\[
\boxed{\text{параллельная архитектура}
\rightarrow
\text{декомпозиция}
\rightarrow
\text{MPI/общая память/GPU}
\rightarrow
\text{ускорение и масштабируемость}}.
\]
Обязательно знать $S_p$, $E_p$, смысл закона Амдала и различие сильной и слабой масштабируемости. Важно подчеркнуть: большее число процессоров само по себе не гарантирует более быстрого расчёта, если возрастают коммуникационные и синхронизационные расходы.

\newpage
